\scanpage{18}\mspage{16}
complexe de De Rham à p.d. \inserted{(ou \unsure{sinon|bien|même} $D^{\cdot\cdot}(X_*,\mathbb{Z})$)} suffit déjà pour récupérer toute la structure mult. \& p.d. coh. de $X_*$ \inserted{(sous réserve de conditions de finitude bien sûr)}. Dans le cas contraire, ce serait cette dernière qui serait le candidat \inserted{alg. ``linéaire''} naturel pour exprimer le type d'homotopie $X_*$ (du moins si $X_*$ \inserted{connexe et} simplement connexe, et en passant à une cat. dérivée convenable bien sûr).
% archive sheet number: 17

